\chapter{Operator Convexity and Jensen Inequalities}\label{ch:operator_convexity}

This chapter collects the operator convexity/concavity results, trace
inequalities, the operator Jensen inequality for positive maps, and the
Lieb concavity theorem. These are standard results in matrix analysis
(Bhatia's \emph{Matrix Analysis}, Chapter~V; \cite[Theorem~5.1]{Wolf2012Quantum};
Hansen--Pedersen's Jensen inequality; and Lieb's concavity theorem). They are
provided here as cited inputs, since the corresponding matrix-analysis proofs for
real powers $x\mapsto x^p$ on the Loewner order and for the matrix logarithm
are not developed elsewhere in this blueprint.

\section{Diagonal Jensen inequality}

\begin{lemma}[Diagonal Jensen inequality]\label{lem:diagonal_jensen_of_convexOn}
    \lean{Matrix.diagonal_jensen_of_convexOn}
    \leanok
    Let $f:\R\to\R$ be convex on $[0,\infty)$, let $A$ be a positive
    semidefinite matrix indexed by a finite type $n$, and let $v\in\C^n$
    satisfy $\langle v,v\rangle=1$. Then
    \[
        f\!(\Re\langle v, A v\rangle)
        \le
        \Re\langle v, f(A)\,v\rangle,
    \]
    where $f(A)$ is defined by the Hermitian continuous functional calculus.
\end{lemma}

\begin{proof}\leanok
    Write $A=U\,\operatorname{diag}(\mu)\,U^{\ast}$ by the spectral
    theorem and set $w=U^{\ast}v$. Since $U$ is unitary,
    $p_i:=\lvert w_i\rvert^2$ satisfies $\sum_i p_i=\lVert v\rVert^2=1$,
    so $(p_i)$ is a probability distribution on $n$. The eigenvalues
    $\mu_i$ lie in $[0,\infty)$ because $A$ is positive semidefinite.
    A direct computation yields
    $\Re\langle v,Av\rangle=\sum_i p_i\mu_i$ and
    $\Re\langle v,f(A)v\rangle=\sum_i p_i f(\mu_i)$, so the scalar Jensen
    inequality applied to the weights $(p_i)$ and points $(\mu_i)$ gives
    the conclusion. See Bhatia, \emph{Matrix Analysis}, Chapter~V.
\end{proof}

\section{Trace concavity and convexity of matrix powers}

\begin{theorem}[Trace concavity of real powers]\label{thm:trace_rpow_concave}
    \lean{TNLean.trace_rpow_concave, trace_rpow_concave}
    \leanok
    \uses{lem:diagonal_jensen_of_convexOn}
    For $p\in[0,1]$, PSD matrices $A_1,A_2$, and $t\in[0,1]$:
    \[
        t\,\Re\tr(A_1^p)+(1-t) \Re\tr(A_2^p)
        \le
        \Re\tr\!((tA_1+(1-t)A_2)^p).
    \]
    Follows from operator concavity of $x\mapsto x^p$ composed with trace
    monotonicity on the Loewner order.
\end{theorem}

\begin{proof}\leanok
    Let $A=tA_1+(1-t)A_2$ and let $\{\psi_j\}$ with eigenvalues
    $\mu_j\ge 0$ be an orthonormal eigenbasis of $A$. A direct
    computation gives $\Re\tr(A^p)=\sum_j \mu_j^p$ and
    $\mu_j=t\,a_j+(1-t)\,b_j$ where
    $a_j=\Re\langle\psi_j,A_1\psi_j\rangle$,
    $b_j=\Re\langle\psi_j,A_2\psi_j\rangle$.
    Scalar concavity of $x\mapsto x^p$ on $[0,\infty)$ yields
    $\mu_j^p\ge t\,a_j^p+(1-t)\,b_j^p$, and the diagonal Jensen
    inequality (Lemma~\ref{lem:diagonal_jensen_of_convexOn},
    concave form) gives
    $a_j^p\le\Re\langle\psi_j,A_1^p\psi_j\rangle$ (and similarly for
    $b_j$). Summing over $j$ and using
    $\sum_j\Re\langle\psi_j,M\psi_j\rangle=\Re\tr M$ yields the
    conclusion. See Bhatia, \emph{Matrix Analysis}, Chapter~V.
\end{proof}

\begin{theorem}[Trace convexity of real powers]\label{thm:trace_rpow_convex}
    \lean{TNLean.trace_rpow_convex, trace_rpow_convex}
    \leanok
    \uses{lem:diagonal_jensen_of_convexOn, thm:trace_rpow_concave}
    For $p\in[1,2]$, PSD matrices $A_1,A_2$, and $t\in[0,1]$:
    \[
        \Re\tr\!((tA_1+(1-t)A_2)^p)
        \le
        t\,\Re\tr(A_1^p)+(1-t) \Re\tr(A_2^p).
    \]
    Follows from operator convexity of $x\mapsto x^p$ composed with trace
    monotonicity on the Loewner order.
\end{theorem}

\begin{proof}\leanok
    Identical eigenbasis reduction as in
    Theorem~\ref{thm:trace_rpow_concave}, with scalar convexity
    (rather than concavity) of $x\mapsto x^p$ on $[0,\infty)$ for
    $p\ge 1$ and the convex form of the diagonal Jensen inequality.
\end{proof}

\section{Operator Jensen inequality for positive maps}

\begin{definition}[Finite-POVM dilation data]\label{def:operator_jensen_povm_aux}
    \lean{TNLean.OperatorJensen.povmIsometry, TNLean.OperatorJensen.povmDiagonal}
    \leanok
    For a finite family of matrices $\{C_i\}_{i \in \iota}$ and a defect block
    $S$, one forms an isometric dilation whose rows are the adjoints of the
    matrices $C_i$ together with the adjoint of $S$. One also forms the scalar
    block-diagonal matrix whose $\iota$-blocks carry prescribed weights and
    whose defect block carries a single scalar.
\end{definition}

\begin{lemma}[Compression auxiliary for Jensen]\label{lem:operator_jensen_inverse_compression}
    \lean{Matrix.PosDef.inverse_compression_le}
    \leanok
    Let $Y$ be positive definite and let $W$ be an isometry. Then
    \[
        (W^\ast Y W)^{-1} \le W^\ast Y^{-1} W.
    \]
\end{lemma}

\begin{proof}
    \leanok
    Apply the Schur-complement criterion to the block matrix
    $\begin{psmallmatrix} Y^{-1} & W \\ W^\ast & W^\ast Y W \end{psmallmatrix}$.
    The lower-right Schur complement is zero, hence positive semidefinite, and
    compressing the resulting upper-left Schur complement by $W$ gives the
    stated inequality.
\end{proof}

\begin{lemma}[Finite-POVM compression identities]\label{lem:operator_jensen_povm_compression}
    \lean{TNLean.OperatorJensen.povmIsometry_star_mul,
        TNLean.OperatorJensen.povmIsometry_compress_diagonal,
        TNLean.OperatorJensen.povm_sum_add_defect,
        TNLean.OperatorJensen.povmDiagonal_posDef,
        TNLean.OperatorJensen.povmDiagonal_inv,
        TNLean.OperatorJensen.povmIsometry_compress_diagonal_inv}
    \leanok
    If the defect block satisfies
    \[
        S S^\ast = \Id - \sum_{i \in \iota} C_i C_i^\ast,
    \]
    then the dilation is an isometry, compressing the weighted scalar
    block-diagonal matrix gives
    \[
        \sum_{i \in \iota} w_i\, C_i C_i^\ast + t\, S S^\ast,
    \]
    the defect relation rewrites as
    \[
        \sum_{i \in \iota} C_i C_i^\ast + S S^\ast = \Id,
    \]
    and strict positivity of all weights together with $t>0$ implies positive
    definiteness of the scalar block-diagonal matrix. If all weights are nonzero
    and the defect scalar is nonzero, the inverse diagonal has reciprocal
    weights, and compressing this inverse gives
    \[
        \sum_{i \in \iota} w_i^{-1}\, C_i C_i^\ast + t^{-1}\, S S^\ast.
    \]
\end{lemma}

\begin{proof}
    \leanok
    Expand the block matrix multiplication entrywise. The defect identity gives
    the isometry relation after summing the $\iota$-blocks and the defect block,
    and the weighted compression identity follows from the same expansion with
    the scalar diagonal inserted. Positive definiteness uses strict positivity
    of every diagonal entry, while the inverse formula uses the nonvanishing of
    every diagonal entry.
\end{proof}

\begin{lemma}[Finite-POVM resolvent inequality]\label{lem:operator_jensen_povm_resolvent}
    \lean{TNLean.OperatorJensen.povm_resolvent_inv_le}
    \leanok
    \uses{lem:operator_jensen_povm_compression, lem:operator_jensen_inverse_compression}
    Let $w_i\ge 0$ and $t>0$. If the defect block satisfies
    \[
        S S^\ast = \Id - \sum_{i \in \iota} C_i C_i^\ast,
    \]
    then
    \[
        \left(\sum_{i \in \iota} w_i\, C_i C_i^\ast + t\,\Id\right)^{-1}
        \le
        \sum_{i \in \iota} (w_i+t)^{-1}\, C_i C_i^\ast + t^{-1}\, S S^\ast.
    \]
\end{lemma}

\begin{proof}
    \leanok
    \uses{lem:operator_jensen_povm_compression, lem:operator_jensen_inverse_compression}
    Apply the compression inequality to the positive definite scalar
    block-diagonal matrix with entries $w_i+t$ and defect entry $t$. The
    compression identities rewrite the compressed matrix as
    $\sum_i w_i C_i C_i^\ast + t\Id$ and its inverse compression as the displayed
    upper bound.
\end{proof}

\begin{theorem}[Jensen for concave real powers]\label{thm:posMap_rpow_concave_jensen}
    \lean{posMap_rpow_concave_jensen}
    \notready
    \uses{def:positive_map, lem:operator_jensen_povm_resolvent}
    For a positive subunital map $T$, a positive semidefinite matrix~$A$,
    and $p\in[0,1]$:
    \[
        T(A^p) \le (TA)^p.
    \]
    This is \cite[Theorem~5.1]{Wolf2012Quantum} for concave~$f(x)=x^p$.
\end{theorem}

\begin{theorem}[Map form of Jensen for concave real powers]\label{thm:rpow_concave_jensen}
    \lean{IsPositiveMap.rpow_concave_jensen}
    \leanok
    \uses{thm:posMap_rpow_concave_jensen}
    This is the corresponding map-form statement of
    Theorem~\ref{thm:posMap_rpow_concave_jensen}.
\end{theorem}

\begin{proof}\leanok
    Direct application of the axiom.
\end{proof}

\begin{theorem}[Jensen for convex real powers]\label{thm:posMap_rpow_convex_jensen}
    \lean{posMap_rpow_convex_jensen}
    \notready
    \uses{def:positive_map}
    For a positive subunital map $T$, a positive semidefinite matrix~$A$,
    and $p\in[1,2]$:
    \[
        (TA)^p \le T(A^p).
    \]
    This is \cite[Theorem~5.1]{Wolf2012Quantum} for convex~$f(x)=x^p$.
\end{theorem}

\begin{theorem}[Map form of Jensen for convex real powers]\label{thm:rpow_convex_jensen}
    \lean{IsPositiveMap.rpow_convex_jensen}
    \leanok
    \uses{thm:posMap_rpow_convex_jensen}
    This is the corresponding map-form statement of
    Theorem~\ref{thm:posMap_rpow_convex_jensen}.
\end{theorem}

\begin{proof}\leanok
    Direct application of the axiom.
\end{proof}

\begin{theorem}[Jensen for concave log]\label{thm:posMap_log_concave_jensen}
    \lean{posMap_log_concave_jensen}
    \notready
    \uses{def:positive_map}
    For a positive \emph{unital} map $T$ and positive-definite~$A$:
    \[
        T(\log A) \le \log(TA).
    \]
    This is \cite[Theorem~5.1]{Wolf2012Quantum} for~$f=\log$.
    Requires unitality ($T(\Id)=\Id$), not merely subunitality.
\end{theorem}

\begin{theorem}[Map form of Jensen for the concave logarithm]\label{thm:log_concave_jensen}
    \lean{IsPositiveMap.log_concave_jensen}
    \leanok
    \uses{thm:posMap_log_concave_jensen}
    This is the corresponding map-form statement of
    Theorem~\ref{thm:posMap_log_concave_jensen}.
\end{theorem}

\begin{proof}\leanok
    Direct application of the axiom.
\end{proof}

\section{Lieb concavity theorem}

\begin{theorem}[Lieb concavity]\label{thm:lieb_concavity_axiom}
    \lean{lieb_concavity_axiom}
    \notready
    For $s\in[0,1]$, any matrix~$K$, and positive-definite matrices
    $A_1,A_2,B_1,B_2$, the map
    $(A,B)\mapsto\tr(K^\dagger A^s K B^{1-s})$
    is jointly concave:
    \[
        t\,\Re\tr(K^\dagger A_1^s K B_1^{1-s})
        +(1-t) \Re\tr(K^\dagger A_2^s K B_2^{1-s})
        \le
        \Re\tr\!(K^\dagger(tA_1+(1-t)A_2)^s K (tB_1+(1-t)B_2)^{1-s}).
    \]
    See \cite{Lieb1973Convex,Ando1979Concavity}.
\end{theorem}

\begin{theorem}[Map form of Lieb concavity]\label{thm:lieb_concavity}
    \lean{lieb_concavity}
    \leanok
    \uses{thm:lieb_concavity_axiom}
    This is the corresponding map-form statement of
    Theorem~\ref{thm:lieb_concavity_axiom}.
\end{theorem}

\begin{proof}\leanok
    Direct application of the axiom.
\end{proof}

\begin{corollary}[Lieb concavity, $K=\Id$ case]\label{cor:lieb_concavity_id}
    \lean{lieb_concavity_id}
    \leanok
    \uses{thm:lieb_concavity}
    For $s\in[0,1]$, $t\in[0,1]$, and positive-definite matrices $A_1,A_2,B_1,B_2$, the
    map $(A,B)\mapsto\Re\tr(A^s B^{1-s})$ is jointly concave:
    \[
        t\,\Re\tr(A_1^s B_1^{1-s})+(1-t) \Re\tr(A_2^s B_2^{1-s})
        \le
        \Re\tr\!((tA_1+(1-t)A_2)^s (tB_1+(1-t)B_2)^{1-s}).
    \]
    Obtained from Theorem~\ref{thm:lieb_concavity} by taking $K=\Id$.
\end{corollary}

\begin{proof}\leanok
    Specialize Theorem~\ref{thm:lieb_concavity} at $K=\Id$ and simplify
    $\Id^{\dagger}=\Id$ together with $\Id\cdot X=X\cdot\Id=X$.
\end{proof}

\begin{corollary}[Concavity of Lieb's map in the first argument]
    \label{cor:lieb_concavity_in_fst}
    \lean{lieb_concavity_in_fst}
    \leanok
    \uses{thm:lieb_concavity}
    For $s\in[0,1]$, $t\in[0,1]$, matrix $K$, and positive-definite matrices $A_1,A_2,B$,
    the map $A\mapsto\Re\tr(K^\dagger A^s K B^{1-s})$ is concave:
    \[
        t\,\Re\tr(K^\dagger A_1^s K B^{1-s})
        +(1-t) \Re\tr(K^\dagger A_2^s K B^{1-s})
        \le
        \Re\tr\!(K^\dagger(tA_1+(1-t)A_2)^s K B^{1-s}).
    \]
\end{corollary}

\begin{proof}\leanok
    Specialize Theorem~\ref{thm:lieb_concavity} at $B_1=B_2=B$, so that
    the convex combination collapses: $tB+(1-t)B=B$.
\end{proof}

\begin{corollary}[Concavity of Lieb's map in the second argument]
    \label{cor:lieb_concavity_in_snd}
    \lean{lieb_concavity_in_snd}
    \leanok
    \uses{thm:lieb_concavity}
    For $s\in[0,1]$, $t\in[0,1]$, matrix $K$, and positive-definite matrices $A,B_1,B_2$,
    the map $B\mapsto\Re\tr(K^\dagger A^s K B^{1-s})$ is concave:
    \[
        t\,\Re\tr(K^\dagger A^s K B_1^{1-s})
        +(1-t) \Re\tr(K^\dagger A^s K B_2^{1-s})
        \le
        \Re\tr\!(K^\dagger A^s K (tB_1+(1-t)B_2)^{1-s}).
    \]
\end{corollary}

\begin{proof}\leanok
    Specialize Theorem~\ref{thm:lieb_concavity} at $A_1=A_2=A$, so that
    the convex combination collapses: $tA+(1-t)A=A$.
\end{proof}